\documentclass[a4paper,12pt]{article}

\usepackage{alltt, fancyvrb, url}
\usepackage{graphicx}
\usepackage[utf8]{inputenc}
\usepackage{float}
\usepackage{hyperref}
\usepackage[italian]{babel}
\usepackage[table]{xcolor}
\usepackage{array}
\usepackage{multirow}
\usepackage{titlesec}
\usepackage{listings}
\usepackage[italian]{cleveref}
\usepackage{acronym}

\sloppy
\setlength{\parindent}{0pt}

\title{\textbf{Joint Project Planning Sessions}}
\author{}
\date{}

\begin{document}
\maketitle

Partecipanti: Project Manager del fornitore e del cliente (ChargeAndTrack), team del cliente, membri principali del team di sviluppo.

\section*{Agenda sessione 1}
\begin{itemize}
    \item Introduzione dello sponsor;
    \item panoramica dello sponsor sul progetto e sulla sua importanza per l'azienda;
    \item introduzione dei co-project manager del fornitore e del cliente;
    \item introduzione dei membri principali del team di sviluppo;
    \item validazione e prioritizzazione dei requisiti;
    \item panoramica sull'approccio che si vuole utilizzare per la pianificazione del progetto;
    \item generazione della WBS e criteri di accettazione;
    \item aggiornamento alla prossima sessione.
\end{itemize}

\section*{Agenda sessione 2}
\begin{itemize}
    \item Introduzione e scopo dela sessione;
    \item discussione e approvazione della WBS e dei criteri di accettazione prodotti durante la prima sessione;
    \item stima della quantità di lavoro, della durata e delle risorse richieste;
    \item aggiornamento alla prossima sessione.
\end{itemize}

\section*{Agenda sessione 3}
\begin{itemize}
    \item Introduzione e scopo dela sessione;
    \item creazione del diagramma delle dipendenze;
    \item individuazione e discussione del critical path e delle date previste per il completamento del progetto e
    delle diverse milestone;
    \item aggiornamento alla prossima sessione.
\end{itemize}

\section*{Agenda sessione 4}
\begin{itemize}
    \item Introduzione e scopo dela sessione;
    \item analisi della schedula prodotta durante la scorsa sessione e sua eventuale compressione;
    \item identificazione dei rischi e degli eventuali piani di mitigazione;
    \item approvazione della schedula finale e del piano di gestione dei rischi.
\end{itemize}

\end{document}